\documentclass[aspectratio=169,11pt]{beamer}
\usepackage[T1]{fontenc}
\usepackage{lmodern}
\usepackage{graphicx}
\usepackage{booktabs}
\usepackage{array}
\newcommand{\DPTitle}{MCP-referenced timing calibration of CMS BTL}
\newcommand{\DPShortTitle}{MCP-referenced timing calibration}
\newcommand{\DPContact}{cms-dpg-conveners-mtd@cern.ch}
% Original Note1.pdf abstract; the slide abstract below remains independent.
\newcommand{\DPCoverAbstract}{%
Several analyses of the 2025 data set benefit from an external time reference
($t_0$). During the testbeam period a Hamamatsu MCP-PMT was placed downstream of
the setup and waveforms were digitized on a Lecroy scope. Since the active area
($\emptyset$=10mm ) of the PMT only covered a couple BTL bars under test and the
scope DAQ captured only 7.5k events per spill, we used a limited data set to
calibrate Reference Bar timing relative to the MCP (which had intrinsic
resolution of 7.5 picoseconds). The calibration functions are provided in this note.}
\newcommand{\DPAbstract}{%
This study examines a timing-calibration procedure for CMS barrel timing layer
(BTL) prototype scintillator bars using September 2025 beam-test data. A microchannel-plate
photomultiplier (MCP-PMT), digitized by an oscilloscope, provides an external
arrival-time reference. A trigger link identifies coincident events in the
BTL readout and the oscilloscope. Segment-dependent clock alignment and
locally weighted smoothing map the reference into the BTL time coordinate.

Quadratic functions of inverse charge describe the charge dependence of the
measured arrival time for four silicon-photomultiplier channels on two reference
bars. These functions provide an empirical time-walk
calibration over the measured charge range. Clock-corrected profiles provide
the reference for the companion dual-threshold study; archived fits are
retained to document the change in calibration procedure. Timing-difference
distributions illustrate the validation strategy. Their interpretation separates
intrinsic reference-detector resolution from residual synchronization effects.
This draft presents the results and the selection, uncertainty and
reproducibility checks that remain before publication.}

\newcommand{\AnalysisRepo}{https://github.com/De-Cristo/LecroyTranslator4CMSTimingDetector}
\newcommand{\AnalysisRevision}{45a907c422d6c8afbe0cdf44cd2827d3433d0f97}
\newcommand{\RepoFile}[2]{\href{\AnalysisRepo/blob/\AnalysisRevision/#1}{#2}}
\usetheme{default}
\setbeamercolor{normal text}{fg=black,bg=white}
\setbeamercolor{structure}{fg=black}
\setbeamercolor{frametitle}{fg=black,bg=white}
\setbeamerfont{frametitle}{size=\large,series=\bfseries}
\setbeamersize{text margin left=7mm,text margin right=7mm}
\setbeamertemplate{navigation symbols}{}
\setbeamertemplate{headline}{%
  \hbox to \paperwidth{\hspace{7mm}\scriptsize\color{black!55}Internal review draft -- 22 October 2026
  \hfill\raisebox{0pt}[9mm][1mm]{\includegraphics[height=7mm]{CMSlogo_black_label_1024.png}}\hspace{6mm}}}
\setbeamertemplate{footline}{%
  \hbox to \paperwidth{\hspace{7mm}\scriptsize\DPShortTitle
  \hfill Draft v1\quad\insertframenumber/\inserttotalframenumber\hspace{7mm}}\vspace{2mm}}
\setbeamertemplate{frametitle}{\vspace{-1mm}\insertframetitle\par\vspace{1mm}}
\setlength{\parskip}{2.5pt}
\newcommand{\Review}[2]{\par\smallskip{\fontsize{7.5}{8.5}\selectfont\color{red!65!black}\textbf{Review #1:} #2\par}}
\newcommand{\Caption}{\fontsize{9}{10.5}\selectfont}
\newcommand{\FigureSlide}[4]{%
  \begin{frame}[t]{#1}
    \begin{columns}[T,onlytextwidth]
      \begin{column}{.61\textwidth}
        \centering\includegraphics[width=\linewidth,height=52mm,keepaspectratio]{figures/#2.pdf}
      \end{column}
      \begin{column}{.37\textwidth}\Caption #3\par\medskip{\fontsize{7}{8.5}\selectfont\color{black!55}#4\par}\end{column}
    \end{columns}
  \end{frame}}
\newcommand{\OpeningSlides}{%
  \begin{frame}
    \centering{\LARGE\DPTitle\par}\vspace{7mm}
    {\large CMS\par}\vspace{3mm}
    \href{mailto:\DPContact}{\DPContact}
    \vspace{7mm}\par{\small Discussion draft; not approved for public release.}
  \end{frame}
  \begin{frame}{Abstract}\small\DPAbstract\end{frame}}
\hypersetup{pdftitle={\DPTitle},pdfauthor={CMS}}

% Generated from frozen numerical results by build_assets.py.
\newcommand{\TimingCount}{10127}
\newcommand{\DBefore}{55.16}
\newcommand{\DAfter}{54.82}
\newcommand{\DLow}{54.32}
\newcommand{\DHigh}{55.22}
\newcommand{\DSD}{61.26}
\newcommand{\DTail}{1.57}
\newcommand{\RBefore}{77.69}
\newcommand{\RAfter}{64.16}
\newcommand{\RLow}{63.27}
\newcommand{\RHigh}{64.99}
\newcommand{\RSD}{70.33}
\newcommand{\RTail}{1.56}
\newcommand{\RefBinsA}{19}
\newcommand{\RefBinsB}{19}
\newcommand{\RefBinsC}{20}
\newcommand{\RefBinsD}{19}
\newcommand{\ExampleAWC}{18.332}
\newcommand{\ExampleLow}{18.120}
\newcommand{\ExampleHigh}{18.607}
\newcommand{\ExampleRMS}{6.80}
\newcommand{\GapCount}{703}
\newcommand{\SummaryChi}{538.6}
\newcommand{\PredictorRMSE}{3.22}
\newcommand{\PredictorFraction}{14.1}

\begin{document}
\OpeningSlides

\FigureSlide{Beam-test setup and the reference signal}{setup}{
The September 2025 beam test used scintillator bars read out at both ends by
silicon photomultipliers (SiPMs). A downstream MCP-PMT provides a fast reference
for a restricted region of the bars.

The BTL front-end readout and oscilloscope use independent time coordinates.
The scope trigger reaches TOFHIR channel 192; a separate 160\,MHz clock link
can also be recorded by the scope.

The reference-bar calibration uses the highlighted lower readout path.
}{Original setup illustration. Detector and readout context:
Ref.~\cite{MTDTDR}; the run-specific setup is described in this study.}

\FigureSlide{Reconstructing the MCP pulse time}{mcp_waveform}{
The figure shows a Hamamatsu MCP waveform and its fitted pulse template.
The fit estimates the pulse centroid rather than using the threshold-crossing
time alone, reducing sensitivity to pulse amplitude.

The scope-trigger time identifies the event; the fitted peak supplies its
relative reference time. Their difference is measured within the same
oscilloscope waveform.
}{}

\FigureSlide{Intrinsic reference timing is a separate measurement}{mcp_pair_timing}{
The distribution shows the difference between the fitted times of two MCP-PMTs
recorded by the oscilloscope. Its measured width is 10.6\,ps.

Assuming independent errors and a 7.5\,ps PHOTONIS contribution gives
$\sigma_{\rm HPK}\simeq\sqrt{10.6^2-7.5^2}=7.5$\,ps.

This scope-only estimate does not include synchronization to the BTL readout.
}{}

\FigureSlide{Matching the same particles between readouts}{event_matching}{
The matching procedure compares consecutive trigger intervals measured by
channel 192 with those recorded by the scope. Correlated intervals identify
candidate matches despite the independent clocks.

The interval-ratio distribution tests matching consistency. It is not a
measurement of the picosecond resolution of the reconstructed particle time.

Only uniquely matched events enter the reference calibration.
}{\RepoFile{FEBD_Sync/simple_sync.py}{GitHub: trigger-matching implementation}.\par
\texttt{dedup\_triggers}: duplicate removal.\par
\texttt{match\_cumulative\_time}: matching and tolerance.}

\FigureSlide{Clock-rate alignment within acquisition segments}{clock_relation}{
The figure compares matched scope-trigger times with TOFHIR trigger times.
A nearly linear relation describes the overall clock-rate difference within
an acquisition segment.

Both an offset and a slope are needed. Even a slope close to one matters
over a long time interval. Clock resets require separate segment coordinates.
}{Archived clock-relation plot. It illustrates synchronization, not
charge-dependent time walk.}

\FigureSlide{Correcting smooth residual clock drift}{clock_residuals}{
After the linear clock-rate correction, the time residual varies
within each of three segments. Points represent matched events; red curves show
the smooth trend estimated by LOWESS, a locally weighted regression~\cite{LOWESS}.

Adding the fitted trend to the linear prediction reconstructs the expected
channel-192 trigger time. The reference difference subtracts this expected
time once from the SiPM time.
}{Archived residuals in ns versus segment time in seconds. These are
pre-LOWESS residuals, not final timing resolutions.
\RepoFile{TimeCalibration/ch192_vs_trigger_lowess.py\#L657-L668}{GitHub: clock prediction};
\RepoFile{TimeCalibration/ch192_vs_trigger_lowess.py\#L816}{reference subtraction}.}

\begin{frame}{A single, explicit time convention}
\small
Let $g(u)$ map a scope-trigger timestamp $u$ into the TOFHIR time coordinate,
including the linear clock relation and the fitted smooth drift.
Define the within-waveform difference
\[
  \delta_{\rm MCP}=t_{\rm peak}^{\rm scope}-t_{\rm trigger}^{\rm scope}.
\]
The reference difference implemented in the calibration is
\[
 \Delta t_i=t_i^{\rm TOFHIR}-g(t_{\rm trigger}^{\rm scope})
                    -\delta_{\rm MCP},
\]
up to a fixed cable and geometry offset. The clock-rate correction for the
short within-waveform interval is negligible and is omitted.

The charge-dependent component of $\Delta t_i$ is the time walk.
An arbitrary constant does not change a shape calibration.\par
{\scriptsize\textbf{GitHub implementation:}
\RepoFile{TimeCalibration/ch192_vs_trigger_lowess.py\#L816}{per-file subtraction};
\RepoFile{TimeCalibration/ch192_vs_trigger_lowess.py\#L1102}{combined-run subtraction}.
Both per-file and combined-run calculations subtract
\texttt{(mcp\_peak - mcp\_trig)}.\par}
\end{frame}

\begin{frame}{Reference bars and calibration selection}
\fontsize{10}{11.5}\selectfont
\begin{tabular}{lll}
\toprule
Upstream bar & SiPM channels & Role\\
\midrule
7 & 133, 154 & First reference pair\\
8 & 137, 150 & Second reference pair\\
\bottomrule
\end{tabular}\par
\medskip
\begin{itemize}
\item Select coincident SiPM hits and a unique MCP-associated event, using
channel 192 as the TOFHIR indicator.
\item Require a successfully reconstructed MCP pulse and consistent
peak-to-trigger timing. This is a waveform-quality cut, not a particle
transit-time measurement inside the MCP.
\item Use $Q$, the TOFHIR ADC charge proxy. It is not a deposited energy in MeV.
\item Fit the profile of $\Delta t_i$ versus $1/Q$ over the retained charge range.
\end{itemize}
{\scriptsize\textbf{GitHub implementation:}\par
\RepoFile{TimeCalibration/ch192_vs_trigger_lowess.py}{Coincidence and walk-sample selection}.\par
\RepoFile{TimeCalibration/ch192_vs_trigger.py}{MCP mapping and optional cuts}.\par
\RepoFile{TimeCalibration/bar_processing.py}{Peak-to-trigger timing selector}.\par}
\end{frame}

\FigureSlide{Archived reference: bar 7, channel 133}{reference_ch133}{
Run 4399 events are grouped in inverse-charge bins. Points show mean time
differences with standard errors of the mean (SEM). The curve is the archived fit,
$W(Q)=a+b/Q+c/Q^2$.

There are 26 retained profile points. Times include the historical display offset.
These archived means use measured channel-192 timing. An unweighted fit
to the means reproduces the archived coefficients; SEM bars are not fit weights.
}{Archived profile and coefficients, displayed without refitting.
The display convention includes a +26\,ns shift.}

\FigureSlide{Archived reference: bar 7, channel 154}{reference_ch154}{
The opposite end of bar 7 is calibrated independently. Points show the mean
time difference and SEM in 20 retained inverse-charge bins.

The red curve is the archived quadratic in $1/Q$, not a quadratic in $Q$.
Channel-dependent coefficients allow for different pulse responses at
the two bar ends.
The archived fit is unweighted and uses measured channel-192 timing.
The retained bin counts are channel dependent.
}{Run 4399 archived calibration, redrawn without refitting.
Same +26\,ns display convention as channel 133.}

\FigureSlide{Archived reference: bar 8, channel 137}{reference_ch137}{
The channel-137 time difference increases as inverse charge increases:
smaller signals cross a fixed threshold later.

Points are profile means with SEM; the red line is the archived quadratic
fit to the profile. There are 48 retained points. The correction applies
within the measured charge range; extrapolation is not validated.
This unweighted result uses measured channel-192 timing and is retained for comparison
with the clock-corrected reference shown next.
}{Run 4399 archived calibration. Coefficients and profile points are
included with the draft.}

\FigureSlide{Archived reference: bar 8, channel 150}{reference_ch150}{
The opposite end of bar 8 has its own charge-dependent delay.
The 25 points show the mean timing difference and SEM; the curve is
the archived quadratic fit to the profile.

The four curves demonstrate the calibration procedure. Each channel retains
its original number of profile points and its original binning.
This is the archived, unweighted fit using measured channel-192 timing.
}{Run 4399 archived calibration. Coefficient units are ns, ns\,ADC and
ns\,ADC$^2$ for $a$, $b$ and $c$.}

\begin{frame}{Clock-corrected reference shared with Note 2}
\footnotesize
The following results use the expected trigger time $g(u)$ reconstructed
from the segment-wise linear clock relation and LOWESS drift:
\[
\Delta t_i=t_i^{\rm TOFHIR}-g(t_{\rm trigger}^{\rm scope})
 -(t_{\rm peak}^{\rm scope}-t_{\rm trigger}^{\rm scope}).
\]
\begin{itemize}
\item Run 4399: exactly five TOFHIR hits, comprising the target upstream pair,
downstream channels 234 and 243, and channel 192, with a unique MCP match.
\item Require finite, positive charge and slope-related quality quantities;
retain $Q$ at or above the left half-maximum boundary of the spectrum peak.
\item Form 20 uniform $1/Q$ bins; retain times within $5(1.4826\,\mathrm{MAD})$
of each bin's median, with at least five entries. Bars show the SEM after clipping.
\item Fit $W(Q)=a+b/Q+c/Q^2$ with an objective weighted by $1/\mathrm{SEM}^2$.
The next four figures use the same reference coefficients as Note 2.
\end{itemize}
{\scriptsize The displayed +26\,ns constant changes neither the shape nor the saved coefficients.}
\end{frame}

\FigureSlide{Clock-corrected reference: bar 7, channel 133}{clock_reference_ch133}{
The time difference now uses the reconstructed trigger time rather than
the measured channel-192 timestamp, suppressing the fitted clock drift.

Points show mean time differences with SEM in \RefBinsA{} retained inverse-charge bins.
The red curve is the SEM-weighted quadratic used in Note 2.

The curve is restricted to retained charge bins and the range between the
outer profile means. Gaps from missing bins are left blank.
}{Run 4399; the reference used in Note 2, redrawn without refitting.
A constant +26\,ns shift is applied only for display.}

\FigureSlide{Clock-corrected reference: bar 7, channel 154}{clock_reference_ch154}{
The same clock reconstruction and profiling procedure calibrates the opposite
end of bar 7. Independent coefficients describe its
charge-dependent timing response.

The \RefBinsB{} points are retained-bin means; vertical bars show SEM after timing
clipping. The red curve is the SEM-weighted quadratic in $1/Q$.
}{The reference used in Note 2. Coefficients and profile points
accompany this note; the +26\,ns shift is applied only for display.}

\FigureSlide{Clock-corrected reference: bar 8, channel 137}{clock_reference_ch137}{
There are \RefBinsC{} retained inverse-charge bins for this channel.
Smaller charge gives a larger timing delay, as expected for fixed-threshold
readout of pulses with different amplitudes.

Points and bars are the mean and SEM after clipping. The saved quadratic
is the MCP benchmark for the channel-137 inverse-slope comparison in Note 2.
}{Run 4399; expected-trigger timing; refitted after clock-quality selection.
The +26\,ns shift affects only the vertical display origin.}

\FigureSlide{Clock-corrected reference: bar 8, channel 150}{clock_reference_ch150}{
The channel-150 reference uses \RefBinsD{} retained profile bins and the same
clock-corrected time definition as the other three channels.

Bars represent SEM, not the spread of individual event times.
The red line is the SEM-weighted quadratic in inverse charge.
These four results connect the MCP calibration in this note to the AWC
shape comparison in Note 2.
}{Run 4399; coefficients are preserved in their original ns convention.
Only the figure receives the constant +26\,ns display shift.}

\begin{frame}{Archived and clock-corrected shapes are close}
\small
The comparison uses both fitted functions in charge intervals populated in both
samples, after removing a constant offset. No multiplicative scaling is fitted.
\begin{center}
\begin{tabular}{lrr}
\toprule
Channel & Shape RMS difference [ps] & Maximum absolute difference [ps]\\
\midrule
133 & 3.20 & 10.63\\
154 & 4.29 & 9.68\\
137 & 1.33 & 3.22\\
150 & 4.91 & 10.47\\
\bottomrule
\end{tabular}

\end{center}
The current reference identifies clock resets from finite timestamps and applies
a clock-only quality requirement. The archived and current versions also differ in trigger
reconstruction, selection, clipping and fit weights.

These are differences between calibration functions, not timing resolutions
or a statistical-equivalence test. Note 2 uses the clock-corrected version.
\end{frame}

\begin{frame}{Clock quality before timing validation}
\small
Segment boundaries are identified from consecutive finite scope timestamps,
including resets across missing entries. Nonfinite trigger times have no valid mapping.

For each topology, file and segment, fit channel 192 against the scope trigger.
Reject the whole group if the line-residual scale $1.4826\,\mathrm{MAD}$ exceeds
1\,ns, before LOWESS. This rule uses clocks only, not SiPM arrival times.

Six bar-8 groups with finite trigger times fail this cut. Changing the threshold
from 1\,ns to 1\,$\mu$s retains the same groups. Suspected one-event association
offsets are excluded; MCP rows are not shifted to repair them.

Refit the channel walk curves after this selection. For the following timing
comparison, both ends must pass the quality and lower-charge cuts. Each end
must also lie within its populated reference bins and outer profile means.
No additional selection is made on either timing-difference observable.

{\scriptsize \RepoFile{Q_Walk/qwalk_t0_data.py}{GitHub: clock-quality selection}.}
\end{frame}

\FigureSlide{Timing validation against the MCP reference}{validated_bar_mcp}{
For \TimingCount{} events, distributions show bar-average minus mapped MCP
time before and after walk correction, with fixed clock mapping.

The core width changes from \RBefore{} to \RAfter\,ps.
The corrected full-sample SD is \RSD\,ps;
\RTail\% of events lie outside the fit window.

The conditional 16--84\% interval is [\RLow, \RHigh]\,ps
(400 file/segment-block bootstraps, fixed calibrations).
Detector, MCP and residual clock effects contribute.
}{Truncated-normal fit: median $\pm2.5(1.4826\,\mathrm{MAD})$.
Median-centered display; full-sample SD; no segment-wise recentering.}

\begin{frame}{What these timing widths establish}
\small
The reproduced distributions compare timing performance before and after
correction under fixed selections and calibrations. They do not measure
intrinsic bar resolution or isolate synchronization jitter.
At fixed position, for end timing errors $e_L,e_R$,
\[
\mathrm{Var}\!\left(\frac{e_L+e_R}{2}\right)
=\frac{\mathrm{Var}(e_L-e_R)}{4}+\mathrm{Cov}(e_L,e_R).
\]
Halving an end-difference width requires controlled position effects and
negligible end-error covariance. Shared MCP and clock errors also prevent
inferring detector covariance directly from the measured reference residuals.

Separately fitted core widths cannot be treated as variances of a common model
to infer synchronization jitter by quadrature subtraction.
This version reports same-event timing validation with conditional statistical
intervals. Intrinsic bar resolution and synchronization-jitter extraction are
outside its scope. Position, covariance, segment-offset and calibration
systematics are not included in the quoted intervals.
\end{frame}

\begin{frame}{Reference-calibration result and scope}
\small
\begin{itemize}
\item A trigger link enables calibration with common events across independent
BTL and oscilloscope time coordinates.
\item A linear clock relation plus a smooth residual correction defines
the reference-time mapping.
\item Four clock-corrected inverse-charge profiles and quadratic functions
provide the reference used in Note 2. The archived fits remain as a comparison;
profile points, SEM and both coefficient sets accompany this draft.
\item Intrinsic MCP timing, reference-bar timing and cross-system
synchronization are different performance quantities.
\item Charge range and operating conditions limit transfer to other
samples; a constant time offset is not a walk-shape difference.
\end{itemize}
\medskip
For the same selected events, walk correction reduces the bar--MCP core width
from \RBefore{} to \RAfter\,ps. This is a conditional timing-validation result,
not an intrinsic-resolution measurement or a complete uncertainty budget.
\end{frame}

% Entries and numbering transcribed from Note1.pdf, printed page 20.
\begin{frame}{References (1/2)}
\fontsize{9}{10.5}\selectfont
\setbeamertemplate{bibliography item}[text]
\begin{thebibliography}{8}
\bibitem[1]{github_repo}
GitHub repository,
\url{https://github.com/De-Cristo/LecroyTranslator4CMSTimingDetector}.
\bibitem[2]{Fabrice}
Presentation by the Saclay group. \url{https://indico.cern.ch/event/1619706/}
\bibitem[3]{picosec}
``Timing Performance of a Micro-Channel-Plate Photomultiplier Tube'',
\url{https://arxiv.org/pdf/1909.12604}
\bibitem[4]{MTDTDR}
CMS Collaboration, ``A MIP Timing Detector for the CMS Phase-2 Upgrade''.
Technical Report CERN-LHCC-2019-003. CMS-TDR-020, CERN, Geneva, 2019 and
``A High-Granularity Timing Detector for the ATLAS Phase-II Upgrade'',
CERN-LHCC-2020-007; ATLAS-TDR-031
\end{thebibliography}
\end{frame}

\begin{frame}{References (2/2)}
\fontsize{9}{10.5}\selectfont
\setbeamertemplate{bibliography item}[text]
\begin{thebibliography}{8}
\bibitem[5]{WhitePresentation}
S. White, Nov. 2025 BTL TB Analysis mtg presentation,
\url{https://indico.cern.ch/event/1614521/contributions/6803723/attachments/3178820/5922530/5vs10vs20GSa_s.pdf}
\bibitem[6]{Leo}
W. R. Leo, ``Techniques For Nuclear and Particle Physics Experiments'',
p.326, Springer-Verlag Berlin Heidelberg.
\bibitem[7]{DualThreshold}
S. White and A. Boletti, ``Amplitude Walk in Fast Timing: The Role of Dual Thresholds''
2026 JINST V21 P02001.
\bibitem[8]{LOWESS}
W. S. Cleveland, ``Robust Locally Weighted Regression and Regression Scatterplots'',
J. Amer. Statist. Assoc. \textbf{74} (1979) 829.
\end{thebibliography}
\end{frame}


\ifdefined\PublicationCandidate\else
\FigureSlide{Backup: clock-phase reconstruction}{clock_edges}{
The reference clock is sampled by the oscilloscope. Several rising-edge
crossings can be reconstructed within a waveform.

Fitting edge time against edge index estimates the clock phase from multiple
cycles. Evaluating the fit near the mean edge index avoids
unnecessary extrapolation.

This is an alternative synchronization diagnostic, distinct from the
trigger-link method used for the archived reference curves.
}{Archived illustration. Correlated edge errors preclude assuming
a $1/\sqrt{N}$ improvement in precision.}

\FigureSlide{Backup: comparing clock estimators}{clock_estimators}{
The recorded MCP time is compared with three clock-phase estimators:
a single edge, a regression intercept and the fitted time at the mean
edge index.

The distributions illustrate the dependence on the estimator. Their widths
include all contributions present in this measurement and do not directly
measure absolute clock accuracy.
The underlying data, common selections and width definitions have not been established.
This internal discussion figure is excluded from the publication candidate.
}{Archived clock-estimator diagnostic, retained for discussion.}
\fi

\end{document}
